\scanpage{511}
\mspage{558}

\[
\diagram{\begin{tikzcd}[ampersand replacement=\&,column sep=large]
L_0 \arrow[r] \& S\deleted{\unsure{\mathcal{T}}}\Gamma_{\rho,\sigma}^{0}\subset B_{\rho,\sigma}\times Z_\rho\times Z_\sigma \\
\ell \arrow[r,mapsto] \& (\ell,1,1)
\end{tikzcd}}\tag{40}
\]
et l'action de $L_0$ par translation à g. sur
$\unsure{S\mathcal{T}'_{\rho,\sigma}|S\Gamma'_{\rho,\sigma}}$
devient, par $(39)$, l'action sur $S\unsure{\mathscr{G}|\mathcal{G}}_{\rho,\sigma}$ donnée par
\[
(\alpha,\beta)\longmapsto(\ell\alpha,\ell\beta).\tag{41}
\]

Posons donc
\margin{\unsure{noté} $\unsure{\Gamma|\mathcal{T}}_{\rho,\sigma}^{0}$ ds \S\,48, $\overline{\text{\uline{IX}}}$ $\unsure{\#}$.}
\[
\deleted{\unsure{\mathcal{T}}}\Gamma_{\rho,\sigma}^{0}
\overset{\mathrm{dfn}}{=}
S\deleted{\unsure{\mathcal{T}}}\Gamma_{\rho,\sigma}^{0}/\operatorname{Im}L_0
\simeq Z_\rho\mathbin{\times}_{\gamma_{\rho,\sigma}}Z_\sigma,
\tag{42}
\]
on trouve que $(39)$ induit une bijection canonique
\[
\diagram{\begin{tikzcd}[ampersand replacement=\&,column sep=large]
\deleted{\unsure{\mathcal{T}}}\Gamma_{\rho,\sigma}^{0}\arrow[r,"\sim"]\&L_0\backslash S\unsure{\mathscr{G}|\mathcal{G}}_{\rho,\sigma}
\end{tikzcd}}
\deleted{\unsure{\simeq Z_\rho\times Z_\sigma}}
\tag{43}
\]
\deleted{L'expression $(37)$ de $S_0\mathcal{T}_{\rho,\sigma}$ donne une suite \inserted{exacte},}
\[
\deleted{(44)\quad 1\longrightarrow L_0\times S}
\qquad\deleted{\text{Notons les suites}}
\]
Posant
\[
\left\{\begin{aligned}
\underset{\cdot}{S}Z_\rho&\overset{\mathrm{dfn}}{=}\operatorname{Ker}(Z_\rho\longrightarrow\Gamma_{\rho,\sigma})=\operatorname{Centr}_{\unsure{\mathcal{G}}}(\rho)\\
SZ_{\rho}&=\operatorname{Ker}(Z_{\rho}\mathbin{-}\Gamma_{\rho,\unsure{2|\sigma}})=\operatorname{Centr}_{\mathcal{G}}(\sigma)
\end{aligned}\right.\tag{44}
\]
\margin{introduire aussi
$S\Gamma_{\rho,\sigma}=\operatorname{Ker}(\Gamma_{\rho,\sigma}\longrightarrow\gamma_{\rho,\sigma}),$
$SG_{\rho,\sigma}=\operatorname{Ker}(G_{\rho,\sigma}\longrightarrow\gamma_{\rho,\sigma})=\mathcal{G}_{\rho,\sigma}=\mathcal{G}$.}

on a donc des suites exactes
\[
\left\{\begin{gathered}
\diagram{\begin{tikzcd}[ampersand replacement=\&,column sep=small]
1\arrow[r,no head]\&SZ_\rho\arrow[r]\&Z_\rho\arrow[r,no head]\&\Gamma_{\unsure{\rho\sigma|\rho,\sigma}}\arrow[r,no head]\&1
\end{tikzcd}}\\
\diagram{\begin{tikzcd}[ampersand replacement=\&,column sep=small]
1\arrow[r,no head]\&SZ_\sigma\arrow[r,no head]\&Z_\sigma\arrow[r,no head]\&\Gamma_{\rho,\sigma}\arrow[r,no head]\&1
\end{tikzcd}}
\end{gathered}\right.\,,\tag{45}
\]
et les expressions $(37)$, $(42)$ de $S\deleted{\unsure{\mathcal{T}}}\Gamma_{\rho,\sigma}^{0}$, $\deleted{\unsure{\mathcal{T}}}\Gamma_{\rho,\sigma}^{0}$ donnent
\[
\diagram{\begin{tikzcd}[ampersand replacement=\&,column sep=small]
1\arrow[r]\&L_0\times SZ_\rho\times SZ_\sigma\arrow[r,no head]\&S_0\deleted{\unsure{\mathcal{T}}}\Gamma_{\rho,\sigma}^{0}\arrow[r,no head]\&\gamma_{\rho,\sigma}\arrow[r,no head]\&1
\end{tikzcd}}\tag{46}
\]
\[
\diagram{\begin{tikzcd}[ampersand replacement=\&,column sep=small]
1\arrow[r]\&SZ_\rho\times SZ_\sigma\arrow[r]\&\deleted{\unsure{S}}\deleted{\unsure{\mathcal{T}}}\Gamma_{\rho,\sigma}^{0}\arrow[r,no head]\&\gamma_{\rho,\sigma}\arrow[r]\&1
\end{tikzcd}}.\tag{47}
\]
\uline{NB} Dans le cas $\mathcal{G}=\mathrm{Sl}(2,\widehat{\mathbb{Z}})$, $\rho$ et $\sigma$ ``modulaires'', on retrouve \deleted{\illegible{1}} les $S_0\deleted{\unsure{\mathcal{T}}}\Gamma_{\rho,\sigma}^{0\unsure{4|\#}}$ et $\deleted{\unsure{\mathcal{T}}}\Gamma_{\rho,\sigma}^{0}$ de \S\,48 $\overline{\text{IX}}$.

\mspage{510}

\scanpage{512}
\mspage{559}

Soit maintenant
\[
\begin{gathered}
S_0M_{\rho,\sigma}^{0*}\subset
S\deleted{\unsure{\mathcal{T}}}\Gamma_{\rho,\sigma}^{0}\times G_{\rho,\sigma}
\quad(\subset B_{\rho,\sigma}\times Z_\rho\times Z_\sigma\times G_{\rho,\sigma})\\
\overset{\mathrm{dfn}}{=}\\
\left\{(u,a,b,U)\,\middle|\,
\underbrace{\delta_B(u)=\delta_\rho(a)=\delta_\sigma(b)=\delta_G(U)}_{\text{relation dans }\gamma_{\rho,\sigma}\ \text{cf }(17)}\right\}
\end{gathered}\tag{49}
\]
\[
\begin{gathered}
M_{\rho,\sigma}^{0}\subset T_\rho\times T_\sigma\times G_{\rho,\sigma}\\
\overset{\mathrm{dfn}}{=}\\
\left\{(a,b,U)\,\middle|\,
\underbrace{\delta_\rho(a)=\delta_\sigma(b)=\delta_G(U)}_{\text{relations dans }\gamma_{\rho,\sigma}}\right\}
\end{gathered}\tag{50}
\]
On a donc un hom canonique
\[
\diagram{\begin{tikzcd}[ampersand replacement=\&,column sep=large]
SM_{\rho,\sigma}^{0}\arrow[r]\&M_{\rho,\sigma}^{0}\\
(u,a,b,U)\arrow[r]\&\deleted{\unsure{(u}}(a,b,U)
\end{tikzcd}}\tag{51}
\]
dont le noyau est formé des $(u,1,1,1)$ où $u\in B_{\rho,\sigma}$ tel que $\delta_B(u)=1$ i.e. $u\in L_0$, d'où une suite exacte
\[
\diagram{\begin{tikzcd}[ampersand replacement=\&,column sep=small]
1\arrow[r,no head]\&L_0\arrow[r]\&SM_{\rho,\sigma}^{0}\arrow[r]\&M_{\rho,\sigma}^{0}\arrow[r,no head]\&1
\end{tikzcd}}\tag{52}
\]
et des hom canoniques
\[
\diagram{\begin{tikzcd}[ampersand replacement=\&,column sep=large,row sep=small]
\&G_{\rho,\sigma}\\
M_{\rho,\sigma}^{0}\arrow[ur,"\Theta_G"]\arrow[r,"\Theta_\rho",very thick]\arrow[dr,"\Theta_\sigma"']\&Z_\rho\\
\&Z_\sigma
\end{tikzcd}}\tag{53}
\]
\[
\left(\text{resp.}\quad
\diagram{\begin{tikzcd}[ampersand replacement=\&,column sep=small,row sep=small]
\&U\\
(a,b,U)\arrow[ur,mapsto]\arrow[r,mapsto]\arrow[dr,mapsto]\&a\\
\&b
\end{tikzcd}}\right),
\]
\unsure{redonnant} les hom analogues sur $SM_{\rho,\sigma}^{0}$ par composition avec $(41)$, et de plus on a un hom canonique
\[
\diagram{\begin{tikzcd}[ampersand replacement=\&,column sep=large]
SM_{\rho,\sigma}^{0}\arrow[r,"\Theta_B"]\&B_{\rho,\sigma}\\
(u,a,b,U)\arrow[r,mapsto]\&u
\end{tikzcd}}.\tag{54}
\]

\mspage{511}

\scanpage{513}

On a aussi \deleted{\unsure{de}} l'interprétation
\[
S_0M_{\rho,\sigma}^{0}\simeq M_{\rho,\sigma}^{0}\mathbin{\times}_{\gamma_{\rho,\sigma}}B_{\rho,\sigma};\tag{55}
\]
( $S_0M_{\rho,\sigma}^{0}$ s'exprime à l'aide de $M_{\rho,\sigma}^{0}$ et $B_{\rho,\sigma}$, et on \unsure{voit} que l'extension de $M_{\rho,\sigma}$ par $L_0$ est l'image inverse de l'extension $B_{\rho,\sigma}$ de $\gamma_{\rho,\sigma}$ par $L_0$. Et on a
\[
M_{\rho,\sigma}^{0}\simeq
\underbrace{Z_\rho\mathbin{\times}_{\gamma_{\rho,\sigma}}Z_\sigma}_{\deleted{\unsure{\mathcal{T}}}\Gamma_{\rho,\sigma}^{0}}
\mathbin{\times}_{\gamma_{\rho,\sigma}}G_{\rho,\sigma}.\tag{56}
\]
d'où on déduit (comme $\delta_\rho,\delta_\sigma,\delta_G$ sont épi)
\[
\diagram{\begin{tikzcd}[ampersand replacement=\&,column sep=small]
1\arrow[r,very thick]\&\deleted{\unsure{\Gamma_\rho|T_\rho}}\underbrace{SZ_\rho\times SZ_\sigma\times\unsure{\mathcal{G}}}_{\overset{\mathrm{dfn}}{=}M_{\rho,\sigma}^{0+}}\arrow[r]\&\unsure{M_{\rho,\sigma}'|M_{\rho,\sigma}^{0}}\arrow[r,"\delta_M^{\unsure{\#|*}}"]\&\gamma_{\rho,\sigma}\arrow[r]\&1
\end{tikzcd}}\tag{57}
\]
où on a \deleted{\unsure{posé}} $(44)$
\[
\deleted{(\unsure{48|58})}\qquad
\left\{\begin{aligned}
SZ_\rho&=Z_\rho\cap\mathcal{G}={\underset{\deleted{\unsure{\mathrm{aut}}}}{\uline{\operatorname{Centr}}}}_{\mathcal{G}}(\rho)\\
SZ_\sigma&=Z_\sigma\cap\mathcal{G}={\underset{\deleted{\unsure{\mathrm{aut}}}}{\uline{\operatorname{Centr}}}}_{\mathcal{G}}(\sigma)
\end{aligned}\right.
\]
De même on a
\[
S_0M_{\rho,\sigma}^{0}\simeq
\underbrace{B_{\rho,\sigma}\mathbin{\times}_{\gamma_{\rho,\sigma}}Z_\rho\mathbin{\times}_{\gamma_{\rho,\sigma}}Z_\sigma}_{S_0\deleted{\unsure{\mathcal{T}}}\Gamma_{\rho,\sigma}^{0}}
\mathbin{\times}_{\gamma_{\rho,\sigma}}G_{\rho,\sigma}
\tag{49}
\]
d'où on a une suite exacte canonique
\[
\diagram{\begin{tikzcd}[ampersand replacement=\&,column sep=small]
1\arrow[r]\&L_0\times SZ_\rho\times SZ_\sigma\times\mathcal{G}\arrow[r]\&S_0M_{\rho,\sigma}^{0}\arrow[r]\&\Gamma_{\rho,\sigma}\arrow[r,no head]\&1
\end{tikzcd}}\tag{50}
\]
Considérons l'hom. can
\[
\diagram{\begin{tikzcd}[ampersand replacement=\&,column sep=large]
\mathcal{G}\arrow[r,"i_{\mathcal{G}}"]\&M_{\rho,\sigma}^{0}\\
U\arrow[r,mapsto]\&(1,1,U)
\end{tikzcd}},\tag{51}
\]

\mspage{512}

\scanpage{514}
\mspage{560}

son image est le troisième facteur dans le noyau de $(57)$, et on a une suite exacte canonique (déduite de $(56)$--$(57)$)
\[
\diagram{\begin{tikzcd}[ampersand replacement=\&,column sep=small]
1\arrow[r]\&\underset{\inserted{\substack{\text{parfois noté}\\\mathcal{G}_{\rho,\sigma}?}}}{\unsure{\mathcal{G}}}\arrow[r]\&M_{\rho,\sigma}^{0}\arrow[r]\&\deleted{\unsure{\mathcal{T}}}\Gamma_{\rho,\sigma}^{0}\arrow[r,no head]\&1
\end{tikzcd}}\tag{52}
\]
\margin{$(53')$ \diagram{\begin{tikzcd}[ampersand replacement=\&,column sep=small]
1\arrow[r]\&SZ_\rho\times SZ_\sigma\arrow[r]\&M_{\rho,\sigma}^{0}\arrow[r]\&\unsure{G|G_1}\arrow[r,no head]\&1
\end{tikzcd}}}
et de même on a (via $(49)$ ou $50)$)
\[
\left\{\begin{aligned}
&(53)\quad\diagram{\begin{tikzcd}[ampersand replacement=\&,column sep=small]
1\arrow[r]\&\mathcal{G}\arrow[r]\&S_0M_{\rho,\sigma}^{0}\arrow[r]\&S_0\deleted{\unsure{\mathcal{T}}}\Gamma_{\rho,\sigma}^{0}\arrow[r]\&1
\end{tikzcd}}.\\[1ex]
&\text{Considérons l'application}\\[1ex]
&(54)\quad\diagram{\begin{tikzcd}[ampersand replacement=\&,column sep=large]
S_0M_{\rho,\sigma}^{0}\arrow[r]\&\mathcal{G}\times\mathcal{G}\times\mathcal{G}\\
\{u,a,b,U\}\arrow[r,mapsto]\&(ua^{-1},ub^{-1},uU^{-1})=(\alpha,\beta,f)
\end{tikzcd}}
\end{aligned}\right.
\]
on \unsure{voit} \deleted{\illegible{1}} par le corollaire p.\,557 (cf $(37)$) que cela induit une bijection
\[
S_0M_{\rho,\sigma}^{0}\simeq(S\unsure{\mathscr{G}|\mathcal{G}}_{\rho,\sigma})\times\mathcal{G}.
\tag{55}
\]
On a ainsi, pour un $x\in S_0M_{\rho,\sigma}^{0}$, d'une part les quantités
\[
\diagram{\begin{array}{cccc}
u,&a,&b,&U\\
\rotatebox{90}{$\in$}&\rotatebox{90}{$\in$}&\rotatebox{90}{$\in$}&\rotatebox{90}{$\in$}\\
B_{\rho,\sigma}&Z_\rho&Z_\sigma&G_{\rho,\sigma}
\end{array}},\tag{56}
\]
associés, d'autre part les éléments
\[
\alpha,\beta,f,\quad\unsure{\zeta|\xi},\eta,\quad\alpha_0\ \beta_0,g,h
\qquad\text{tous }\in\mathcal{G}\tag{57}
\]
\inserted{où $(\alpha,\beta,f)$} définis en termes de $(56)$ dans $(54)$, et où on a posé
\[
\left\{\begin{aligned}
\zeta&=\alpha\rho(\alpha)^{-1},\quad
\eta=\beta\sigma(\beta)^{-1},\quad
\alpha_0=f^{-1}\alpha\unsure{\#|\zeta},\quad\beta_0=f^{-1}\beta,\\
g&=\alpha_0\rho(\alpha_0)^{-1}=f^{-1}\zeta\rho(f),\quad
h=\beta_0\sigma(\beta_0)^{-1}=f^{-1}\eta\sigma(f)
\end{aligned}\right.\tag{58}
\]

\mspage{513}

\scanpage{515}

On aura, en termes de $\alpha,\beta,f$, les expressions suivantes de l'action de $U$ sur $\mathcal{G}$ :
\[
\left\{\begin{aligned}
U(\rho)&=\operatorname{int}(f^{-1}\alpha)\deleted{\unsure{\operatorname{int}(\alpha)}}\rho
 =\operatorname{int}(f^{-1})(\zeta\rho)\\
U(\sigma)&=\operatorname{int}(f^{-1}\beta)\deleted{\unsure{\operatorname{int}(\beta)}}(\sigma)
 =\operatorname{int}(f^{-1})(\eta\sigma).
\end{aligned}\right.\tag{59}
\]
\deleted{\unsure{et}\ $\unsure{\zeta|\xi},\unsure{\eta}$\unsure{.}}

\uline{\unsure{Remarque|Remarq.}} Finalement, dans tous ces développements, on n'a pas vraiment utilisé la soi-disante ``hyp. fondamentale'' $(20)\deleted{\unsure{b}}$, \deleted{\unsure{et}} sauf la partie a). En l'absence de b), il y a lieu d'introduire
\[
\begin{gathered}
\Sigma_{\rho,\sigma}^{\mathrm{eff}}\subset\Sigma_{\rho,\sigma}\subset\mathcal{G}\times\mathcal{G}\\[-1pt]
\begin{array}{c@{\qquad}l}
\rotatebox{90}{$=$}&\\[-3pt]
\multicolumn{2}{l}{\left\{(\zeta,\eta)\in\unsure{\mathcal{G}|\Sigma_\rho}\times\mathcal{G}\mid
\exists u\in B_{\rho,\sigma},\ u(\rho)=\zeta\rho,\ u(\sigma)=\eta\sigma\right\},}
\end{array}
\end{gathered}\tag{60}
\]
qui n'est donc autre que l'image de l'application $(9)$ : couples \uline{effectifs} $(\zeta,\eta)$ satisfaisant les conditions définissant \deleted{\unsure{les}\ $\unsure{(5)}$}\inserted{$\Sigma_{\rho,\sigma}$ (cf $(8)$)}\unsure{)}

\deleted{L'application
$\diagram{\begin{tikzcd}[ampersand replacement=\&,column sep=large]
S\unsure{\mathscr{G}|\mathcal{G}}_{\rho,\sigma}\arrow[r]\&\Sigma_{\rho,\sigma}
\end{tikzcd}}$
\inserted{$(28)$} (ne sera pas néc. surjective, son image \unsure{est|sera} $\Sigma_{\rho,\sigma}^{\mathrm{eff}}\inserted{\simeq B_{\rho,\sigma}}$), d'ailleurs il n'y a rien à changer.}
\margin{\unsure{puis|par}\par\unsure{bij.|bis}}

notons
\[
\begin{aligned}
S\unsure{\mathscr{G}|\mathcal{G}}_{\rho,\sigma}^{\mathrm{eff}}
 &=\text{Image inverse de }\Sigma_{\rho,\sigma}^{\mathrm{eff}}\\
 &\text{par }\diagram{\begin{tikzcd}[ampersand replacement=\&,column sep=large]
S\unsure{\mathscr{G}|\mathcal{G}}_{\rho,\sigma}\arrow[r,no head]\&\Sigma_{\rho,\sigma}
\end{tikzcd}}\quad(28),
\end{aligned}\tag{61}
\]
\deleted{l'image inverse}
et c'est sur $S\unsure{\mathscr{G}|\mathcal{G}}_{\rho,\sigma}^{\mathrm{eff}}$ qu'on aura une structure

\mspage{514}

\scanpage{516}
\mspage{561}

de groupe \unsure{isomorphe} \unsure{à|:} $S_0\mathcal{T}_{\rho,\sigma}^{0}$, par l'isomorphisme (qui remplace (38))
\diagram{\begin{tikzcd}[ampersand replacement=\&,column sep=large]
 S_0\deleted{\unsure{\mathcal{T}}}\Gamma_{\rho,\sigma}^{0} \arrow[r,phantom,"\simeq"] \& S_0\unsure{\mathscr{G}|\mathcal{G}}_{\rho,\sigma}^{\mathrm{eff}}\,. \\
 (\underline{u},\underline{a},\underline{b}) \arrow[r] \& (\underbrace{ua^{-1}}_{\zeta},\underbrace{ub^{-1}}_{\eta})
\end{tikzcd}\quad(62)}

Ceci posé, revenons à $M_{\rho,\sigma}^{\sim}$, et la suite exacte
\diagram{\begin{tikzcd}[ampersand replacement=\&,column sep=small]
 1 \arrow[r] \& \mathcal{G}_{0,3}^{+\wedge} \arrow[r] \& M_{\unsure{0,3|\rho,\sigma}}^{\sim} \arrow[r] \& \underbrace{\Gamma_{0,3}^{\sim}}_{\Gamma_{\mathbb Q}^{\sim}} \arrow[r] \& 1
\end{tikzcd}\quad(63).}

Considérons le sous-groupe
\[
\begin{aligned}
(64)\qquad
\Gamma_{\mathbb Q}^{\prime\sim}
 &= \operatorname{Ker}\bigl(\Gamma_{0,3}^{\sim}
 \xrightarrow{(\varepsilon_3,\deleted{\unsure{\varepsilon_2}})}
 \{\pm1\}\deleted{\unsure{\times\{\pm1\}}}\bigr)\\
 &= \{\gamma\in\Gamma_{0,3}^{\sim}\mid\mu(\gamma)\equiv1\ (\unsure{3}\deleted{\unsure{12}})\}.
\end{aligned}
\]
\deleted{\illegible{3}}
\deleted{\unsure{deux facteurs}}

et soit \deleted{$\unsure{M_{0,3}^{\prime\sim}}$}
\[
 (65)\qquad M_{0,3}^{\prime\sim}
 = \text{\unsure{Image inv.} de }\Gamma_{0,3}^{\prime\sim}\text{ dans }M_{0,3}^{\sim}.
\]

Soit $\mathcal{G}$ un groupe profini, et
\[
 (66)\qquad\psi:\mathcal{G}_{0,3}^{+\wedge}\longrightarrow\mathcal{G},
\]
un hom. \inserted{continu} surjectif --- ce qui revient au $\overline{\mathrm m}$ que de se donner deux éléments
\[
 (67)\qquad\rho_{\mathcal{G}},\unsure{\sigma_{\mathcal{G}}}\in\mathcal{G}
\]
(les images de $\rho,\sigma$ par $\psi$) satisfaisant
\[
 (68)\qquad\rho_{\mathcal{G}}^3=\sigma_{\mathcal{G}}^2=1.
\]
\mspage{515}

\scanpage{517}

\margin{l'introduire par conditions équivalentes :

\unsure{a)}
\diagram{\begin{tikzcd}[ampersand replacement=\&,column sep=large]
 M_{0,3}^{\prime\sim}[0] \arrow[r] \& \Sigma_{\rho,\sigma}^{\text{\textcircled{eff}}}
\end{tikzcd}} \unsure{tout}

b) $\forall u\in M_{0,3}^{\prime\sim},\ \exists u_{\unsure{\mathcal{G}}}\in\operatorname{Aut}\mathcal{G}$ tel que
$u_{\unsure{M|\mathcal{G}}}\circ\psi=\psi\circ u$

b') Idem avec $u\in M_{0,3}^{\prime\sim}[0]$

\unsure{c)} $\exists\psi_M$ qui prolonge $\psi$

c') $\exists\psi_G$ qui prolonge $\psi$

On dira donc que \deleted{\unsure{eff.}} \unsure{automatique|autom.}

$\psi$ est \uline{\unsure{effectif}} si (20) b) satisfait}

\uline{Définition} On dira que $\psi$ est \uline{effectif} si pour tout
$u\in M_{0,3}^{\prime\sim}[0]$, posant $u(\rho)=\zeta\rho$, $u(\sigma)=\eta\sigma$,
\[
 (\zeta_{\mathcal{G}},\eta_{\mathcal{G}})
 =(u(\zeta),u(\eta))\in\Sigma_{\rho,\sigma}
\]
effectif (i.e. $\in\Sigma_{\rho,\sigma}^{\mathrm{eff}}$, i.e. $\exists u_{\mathcal{G}}$ \unsure{autom.} $\mathcal{G}$ tel qu'on ait \deleted{$\unsure{u_{\mathcal{G}}\rho}$}
\[
 (69)\qquad
 u_{\mathcal{G}}(\rho_{\mathcal{G}})=\zeta_{\mathcal{G}}\rho_{\mathcal{G}},
 \qquad
 u_{\mathcal{G}}(\sigma_{\mathcal{G}})=\eta_{\mathcal{G}}\sigma_{\mathcal{G}}
\]
ou, ce qui revient au $\overline{\mathrm m}$, qu'on a commutativité dans
\diagram{\begin{tikzcd}[ampersand replacement=\&,column sep=large,row sep=large]
 \mathcal{G}_{0,3}^{+\wedge} \arrow[r,"u"] \arrow[d,"\psi"'] \& \mathcal{G}_{0,3}^{+\wedge} \arrow[d,"\psi"] \\
 \mathcal{G} \arrow[r,"u_{\mathcal{G}}"'] \& \mathcal{G}
\end{tikzcd}\quad(70).}

(Donc dire que c'est le cas pour \unsure{tt} $u\in M_{0,3}^{\prime\sim}\deleted{[0]}$ signifie que ces $u$ transforment $\operatorname{Ker}\psi$ en lui $\overline{\mathrm m}$ --- ou aussi, que
$M_{0,3}^{\prime\sim}\deleted{\unsure{[0]}}$
$(=M_{0,3}^{\prime\sim}[0].\mathcal{G}_{0,3}^{+\wedge})$ tout entier stabilise $\operatorname{Ker}\psi$.
\deleted{Quand} Cette condition est donc automatiquement remplie quand $\mathcal{G},\rho_{\mathcal{G}},\sigma_{\mathcal{G}}$ satisfont \unsure{la condition|les conditions} b) de (20).

On a fait ce qu'il fallait pour avoir \unsure{le}
\mspage{516}

\scanpage{518}
\mspage{562}

\uline{Théorème} Soit $\psi=\psi_{\mathcal{G}}:\mathcal{G}_{0,3}^{+\wedge}\longrightarrow\mathcal{G}$ un hom surjectif effectif (cf ci-dessus).
Alors il existe un hom unique
\[
 (71)\qquad\psi_M:\deleted{\unsure{S_0}}M_{0,3}^{\prime\sim\deleted{\unsure{*}}}
 \longrightarrow\deleted{\unsure{S_0}}M_{\rho_{\mathcal{G}},\sigma_{\mathcal{G}}}^{0}
\]
\deleted{rendant commutatif le diagramme}
\diagram{\deleted{\begin{tikzcd}[ampersand replacement=\&,column sep=large,row sep=large]
 \mathcal{G}_{0,3}^{+\wedge} \arrow[r,"\unsure{\psi|\varphi}"] \arrow[d,"\psi"'] \& \deleted{\unsure{\mathcal{G}}}\,M_{0,3}^{\prime\sim} \arrow[d,"\psi"] \\
 \deleted{\unsure{\mathcal{G}}}\,\mathcal{G} \arrow[r] \& M_{\rho_{\mathcal{G}},\sigma_{\mathcal{G}}}^{0}\ (\overset{\mathrm{dfn}}{=}M_{\psi}^{0})
\end{tikzcd}\quad(72)}}

\margin{En fait, on a \unsure{bien sûr en} premier lieu un hom
\diagram{\begin{tikzcd}[ampersand replacement=\&,column sep=large]
 S_0M_{0,3}^{\prime\sim} \arrow[r] \& S_0M_{\psi}^{0}.
\end{tikzcd}}}

ayant la propriété suivante : si \deleted{$\unsure{u\in}$}
$u\unsure{=|\in}M_{0,3}^{\prime\sim}$ est de la forme $u_{\alpha,\beta,f}$
($\alpha,\beta,f\in\mathcal{G}_{0,3}^{+\wedge}$, satisfaisant l'équation des lacets
$\alpha\rho(\alpha)^{-1}=\beta\sigma(\beta)^{-1}\ell_1^{\nu}$), alors
$\psi_M(u)$ est l'élément de \deleted{$\unsure{M_{\rho_{\mathcal{G}},\sigma_{\mathcal{G}}}^{0}}$} $M_{\psi}^{0}$ défini par les invariants $\psi(\alpha),\psi(\beta);\psi(f)$.

\uline{Corollaire} On a alors des hom de suites exactes
\diagram{\begin{tikzcd}[ampersand replacement=\&,column sep=small,row sep=large]
 1 \arrow[r] \& \mathcal{G}_{0,3}^{+\wedge} \arrow[r,no head] \arrow[d,"\psi_{\mathcal{G}}"'] \& M_{0,3}^{\prime\sim} \arrow[r] \arrow[d,"\psi_M"] \& \Gamma_{\mathbb Q}^{\prime\sim} \arrow[r] \arrow[d,"\unsure{\psi_{\mathcal{T}}|\psi_{\Gamma}}"] \& 1 \\
 1 \arrow[r] \& \mathcal{G} \arrow[r] \& M_{\psi}^{0} \arrow[r] \& \deleted{\unsure{\mathcal{T}}}\Gamma_{\psi}^{0} \arrow[r] \& 1
\end{tikzcd}\quad(72).}

(L'hom. $\psi_M$ s'explicite à l'aide de trois homomorphismes
\mspage{517}

\scanpage{519}

\[
\left\{\begin{aligned}
\psi_{T_\rho}&:\Gamma_{\mathbb Q}^{\prime\sim}\longrightarrow \unsure{Z_\rho|T_\rho}\\
\psi_{T_\sigma}&:\Gamma_{\mathbb Q}^{\sim}\longrightarrow \unsure{Z_\sigma|T_\sigma}
\end{aligned}\right.\qquad(73)
\]
\[
\psi_G:M_{0,3}^{\prime\sim}\longrightarrow G_{\rho,\sigma}\qquad(74)
\]
\margin{NB on écrit pour simplifier $\rho,\sigma$ au lieu de $\rho_{\mathcal G},\sigma_{\mathcal G}$ en indices}

\deleted{satisfaisant les conditions}
(Pour \deleted{$\unsure{u}$} $x\in M_{0,3}^{\prime\sim}$,
\inserted{$\dot x$ son image dans $\Gamma_{\mathbb Q}^{\prime\sim}$,}
\deleted{$a=\psi_{\unsure{T_\rho|\Gamma_\rho}}(\dot x),\ b=\psi_{\unsure{T_\rho|\Gamma_\rho}}(\unsure{y|\dot x}),\ U=\unsure{\psi u|\psi(x)}$}
\[
\deleted{\unsure{\delta_G(}}\quad
 a=\psi_{\unsure{T_\rho|\Gamma_\rho}}(\dot x),\quad b=\psi_{\unsure{T_\sigma|\Gamma_\sigma}}(\dot x),\quad U=\psi_G(x),
\qquad(75)
\]
on a
\[
\delta_\rho(a)=\delta_\sigma(b)=\delta_G(U)\quad\text{dans }\gamma_{\rho,\sigma}\qquad(76)
\]
i.e. les trois hom composés
\diagram{\begin{tikzcd}[ampersand replacement=\&,column sep=large,row sep=large]
\&\Gamma_{\mathbb Q}^{\prime\sim}\arrow[r,"\psi_{T_\rho}"]\&T_\rho\arrow[rd,"\delta_\rho"]\&\\
M_{0,3}^{\prime\sim}\arrow[ru,"\mathrm{can}"]\arrow[r,"\mathrm{can}"]\arrow[rrd,"\psi_G"',"\unsure{1}"{pos=.43}]\arrow[rrd,shorten >=55pt,shorten <=40pt]\&\Gamma_{\mathbb Q}^{\prime\sim}\arrow[r,no head,"\psi_{T_\sigma}"]\&T_\sigma\arrow[r,"\delta_\sigma"]\&\gamma_{\rho,\sigma}\arrow[ld,"\delta_G"]\\
\&\&G_{\rho,\sigma}\&
\end{tikzcd}\qquad(77)}
sont égaux. De plus, on aura d'après (71),
\[
\psi_G|\mathcal G_{0,3}^{+\wedge}=\psi_{\mathcal G},\qquad(78)
\]
et enfin ceci :

(79) Si $x\in \deleted{\unsure{M}}M_{0,3}^{\prime\sim}$, $U=\psi_G(x)$, on a commutativité dans
\diagram{\begin{tikzcd}[ampersand replacement=\&,column sep=huge,row sep=large]
\mathcal G_{0,3}^{+\wedge}\arrow[r,"\psi"]\arrow[d,"\operatorname{int}(x)|\mathcal G_{0,3}^{+\wedge}"']\&\mathcal G\arrow[d,"\operatorname{int}(U)|\mathcal G"]\\
\mathcal G_{0,3}^{+\wedge}\arrow[r,"\psi"]\&\mathcal G
\end{tikzcd}\qquad(80)}

\hfill $.\ .\ .\ .\ .\ .$ 
\mspage{518}

\scanpage{520}
\mspage{563}

. On peut regarder tout ceci comme un énoncé de \uline{fonctorialité} des constructions faites, p.r. à un hom surjectif
\diagram{\begin{tikzcd}[ampersand replacement=\&]
\mathcal G\arrow[r,"\varphi=\varphi_{\mathcal G}"]\&\mathcal G'
\end{tikzcd}\qquad(81)}
\margin{\unsure{Il aurait} mieux valu prendre \unsure{donc}
\diagram{\begin{tikzcd}[ampersand replacement=\&]\mathcal G'\arrow[r,"\varphi"]\&\mathcal G\end{tikzcd}},
pour pouvoir spécialiser en $\mathcal G'$}
transformant le couple $(\rho,\sigma)$ de générateurs topologiques de $\mathcal G$ en un couple $(\rho',\sigma')$.
Il y a alors un unique homomorphisme
\diagram{\begin{tikzcd}[ampersand replacement=\&]
S_0M_{\mathcal G}^{0}\arrow[r,"\varphi_{M^0}"]\&S_0M_{\mathcal G'}^{0}
\end{tikzcd}\qquad(82)}
satisfaisant les conditions suivantes :

\margin{ceci montre que $\varphi_{SM}$ est déterminé quand on connaît les quatre hom $\varphi_B,\varphi_\rho,\varphi_\sigma,\varphi_G$}
a) $\unsure{\varphi_M|\varphi_{SM}}$ \unsure{passe} en quotient \unsure{en|à} des hom
\[
\left\{\begin{aligned}
\varphi_B&:B_{\rho,\sigma}\longrightarrow\unsure{Z_{\rho,\sigma}|\Sigma_{\rho,\sigma}}\\
\varphi_\rho&:Z_\rho\longrightarrow Z_{\rho'}\\
\varphi_\sigma&:Z_\sigma\mathrel{\relbar\joinrel\relbar} Z_{\sigma'}\\
\varphi_G&:G_{\rho,\sigma}\mathrel{\relbar\joinrel\relbar} G_{\rho',\sigma'}
\end{aligned}\right.\ ;\qquad(83)
\]

b) \deleted{On a commutativité dans le diagramme,}
\inserted{\deleted{Les hom $\varphi_B,\varphi_\rho,\varphi_\sigma$ \unsure{sont induits par}}}
\[
(84)\quad\left\{\begin{gathered}
\deleted{\text{\unsure{uniquement}}}\quad\varphi_G\qquad
\text{(ce qui montre que la connaissance}\\
\text{\illegible{1} de }\varphi_{SM}\ \unsure{b|se}\ \text{revient à celle de }\varphi_G\text{)};
\end{gathered}\right.
\]

c) On a commutativité dans
\diagram{\begin{tikzcd}[ampersand replacement=\&,column sep=large,row sep=large]
\mathcal G\arrow[r]\arrow[d,"\varphi=\varphi_{\mathcal G}"']\&G_{\rho,\sigma}\arrow[d,"\varphi_G"]\\
\mathcal G'\arrow[r]\&G_{\rho',\sigma'}
\end{tikzcd}\qquad(85)}
(ce qui montre, puisque $G_{\rho,\sigma}=B_{\rho,\sigma}\cdot\mathcal G$, que $\varphi_G$ est connu quand on connaît $\varphi_B$ ;

d) Pour $u\in B_{\rho,\sigma}$, \deleted{\unsure{désignant} par $u'$} \inserted{posant $u'=\varphi_B(u)$}, on a commutativité dans
\diagram{\begin{tikzcd}[ampersand replacement=\&,column sep=large,row sep=large]
\mathcal G\arrow[r,"u"]\arrow[d,"\varphi_{\mathcal G}"']\&\mathcal G\arrow[d,"\varphi_{\mathcal G}"]\\
\mathcal G'\arrow[r,"u'"]\&\mathcal G'
\end{tikzcd}\qquad(86)}
\mspage{519}

\scanpage{521}

-- (ce qui détermine $\varphi_B$ en termes de $\varphi=\varphi_{\mathcal G}$, donc \unsure{cela|elle} détermine $\unsure{\varphi_{S_0M}}$, et par passage au quotient, aussi
\[
 (87)\qquad \varphi_M:M_{\rho,\sigma}\longrightarrow M_{\rho',\sigma'}.
\]

Une autre façon, plus directe, de décrire $\varphi_{S_0M^0}$ ou $\varphi_M$, en termes des \deleted{coordonnées} \inserted{``paramètres''} $(\alpha,\beta,f)$ au lieu des $(u,a,b,U)$, est de dire qu'on \deleted{$\unsure{\varphi_{S_0M^0}}$ \unsure{transforme}} transforme un $u_{\alpha,\beta,f}$ en le $u_{\alpha',\beta',f'}$, où $\alpha',\beta',f'$ sont les images de $\alpha,\beta,f$ par $\varphi$:
\[
 (88)\qquad
 \varphi_{\deleted{\unsure{\mathcal G}}S_0\unsure{M_0|M^0}}
 \left(\underbrace{u_{\alpha,\beta,f}}_{\in\,S_0M^0_{\rho,\sigma}}\right)
 =u_{\alpha',\beta',f'},\quad
 \left\{\begin{aligned}
 \alpha'&=\varphi(\alpha),\\
 \beta'&=\varphi(\beta),\\
 f'&=\varphi(f).
 \end{aligned}\right.
\]

On retrouve, ainsi, sous une forme un peu plus précise, le \unsure{th.} précédent, en \deleted{\unsure{mettant}} l'appliquant au cas où $\mathcal G=\mathcal G_{0,3}^{+\wedge}$ (le ``cas universel'')$^*$, $\mathcal G'$ quelconque --- il \unsure{vaut} mieux dans la notation (81) interchanger le rôle de $\mathcal G$ et $\mathcal G'$.
\margin{$^*$ tout au moins ``universel'' p.r. à l'hyp. \unsure{suppl.} $\rho^3=1,\ \sigma^2=1$ dans $\mathcal G$.}
Il est bon d'\deleted{\unsure{expliciter}}\inserted{expliciter les invariants essentiels} (dans le cas ``universel'':

\deleted{pour}
\[
\begin{aligned}
 (90)\qquad
 \mathcal G'&=\mathcal G_{0,3}^{+\wedge},\qquad
 B_{\rho',\sigma'}^{\deleted{\unsure{0}}\,\prime}
 =M_{0,3}^{\prime\sim}[0],\qquad
 \gamma_{\rho',\sigma'}^{\unsure{\prime|0}}
 =\Gamma_{\mathbb Q}^{\prime\sim},\\
 G_{\rho',\sigma'}^{\unsure{\prime|0}}
 &=M_{0,3}^{\prime\deleted{\unsure{0}}\sim},\qquad
 Z_{\rho}^{\prime\deleted{\unsure{0}}}
 =M_{0,3}^{\prime\sim}\unsure{\left(-\frac13\right)|\left(\frac13\right)},\qquad
 Z_{\sigma}^{\prime\deleted{\unsure{0}}}
 =M_{0,3}^{\prime\sim}\unsure{\left(\frac12\right)|\left(-\frac12\right)},\\
 \Gamma_{\rho',\sigma'}^{\prime}&=\Gamma_{\mathbb Q}^{\prime\sim}
 \quad\deleted{\left(\text{d'où}\quad
 Z_{\rho}^{\prime}\xrightarrow{\sim}\gamma_{\rho',\sigma'}^{\unsure{\prime|0}},\quad
 Z_{\sigma}^{\prime}\xrightarrow{\sim}\gamma_{\rho',\sigma'}^{\prime}\right)}.
\end{aligned}
\]
\inserted{d'où}
\mspage{520}


\scanpage{522}
\mspage{564}
\[
\begin{gathered}
\Gamma_{\rho',\sigma'}^{\prime}\xrightarrow{\sim}\gamma_{\rho',\sigma'}^{\prime},\quad
Z_{\rho'}^{\prime}\xrightarrow{\sim}\gamma_{\rho',\sigma'}^{\prime},\quad
Z_{\sigma'}^{\prime}\xrightarrow{\sim}\gamma_{\rho',\sigma'}^{\prime},\quad\text{donc}\\
\unsure{(91)|(92)}\qquad
S_0M_{\rho',\sigma'}^{\prime}\simeq
\overset{\inserted{S_0G_{\rho',\sigma'}^{\prime}\unsure{\simeq|=}}}{S_0M_{0,3}^{\prime\sim}},\quad
\deleted{\unsure{G_{\rho',\sigma'}=}}M_{\rho',\sigma'}^{\prime}\xrightarrow{\sim}
\overset{\inserted{G_{\rho',\sigma'}^{\prime}\xrightarrow{\sim}}}{\deleted{\unsure{G}}M_{0,3}^{\prime\sim}}.
\end{gathered}
\]
\margin{d'où
\[
(92)\quad\left\{\begin{aligned}
SZ_{\rho'}^{\prime}&=SZ_{\sigma'}^{\prime}\\
&=S\unsure{\Gamma|\mathcal T}_{\rho',\sigma'}^{\prime}=1
\end{aligned}\right.
\quad(\unsure{cf.})\ \text{de}\ (92)
\]}

\uline{Remarque} En écrivant les relations (86), il devenait évident que les notations utilisées dans cette section II (depuis p.\ \unsure{0}\ 553') étaient inadéquates, et qu'il convient de les modifier de la façon suivante:
\margin{\unsure{Ceci m'avait été inspiré dans le cas} $\mathcal G=\mathrm{Sl}(2,\widehat{\mathbb Z})$, \unsure{couple modulaire} $\rho,\sigma$ \unsure{qui ne résulte pas de} $\mathcal G_{0,3}^{+\wedge}$, à cause des $\{\pm1\}$! \unsure{mais dans ce cas} \illegible{1} \unsure{universel}, aurait dû me servir de fil conducteur pour mes notations.}
\[
\begin{aligned}
(93)\qquad&\Sigma_{\rho,\sigma},\ B_{\rho,\sigma},\
\overset{\inserted{S_0G_{\rho,\sigma},}}{G_{\rho,\sigma}},\
\gamma_{\rho,\sigma},\ \deleted{\unsure{\Gamma_{\rho,\sigma}}},\
Z_\rho,\ Z_\sigma,\ \mathscr G_{\rho,\sigma},\ S\mathscr G_{\rho,\sigma}\\
&Z_\rho,\ Z_\sigma,\ S_0\Gamma_{\rho,\sigma}^{\deleted{\unsure{*}}0},\
\Gamma_{\rho,\sigma}^0,\ SZ_\rho,\ SZ_\sigma,\ S_0M_{\rho,\sigma}^0,\ M_{\rho,\sigma}^0.
\end{aligned}
\]

On met partout un accent $'$ pour en faire $\Sigma'_{\rho,\sigma}$, $B'_{\rho,\sigma}$ etc, et on enlève l'exposant $^0$ de $S_0\Gamma_{\rho,\sigma}^0$, $\Gamma_{\rho,\sigma}^0$, $S_0M_{\rho,\sigma}^0$, $M_{\rho,\sigma}^0$ (qui se trouve donc remplacé par un $'$).

Autres cas particuliers importants à \unsure{expliciter}:

$2^\circ)$ $\mathcal G=\mathrm{Sl}(2,\widehat{\mathbb Z})/\pm1$, $(\rho,\sigma)$ couple modulaire.
\[
\begin{aligned}
(94)\quad\left\{\begin{aligned}
B'_{\rho,\sigma}&\simeq
\left\{\begin{pmatrix}\mu&\lambda\\0&1\end{pmatrix}
\mid\mu\in\widehat{\mathbb Z}^{*},\ \mu\equiv1\ (3),\ \lambda\in\widehat{\mathbb Z}\right\}\\
&\inserted{L_0=\left\{\begin{pmatrix}1&\lambda\\0&1\end{pmatrix}\mid\lambda\in\widehat{\mathbb Z}\right\}}\\
&\deleted{\unsure{\mathrm{Gl}(\widehat{\mathbb Z})/\mu_2(\mathbb Z)}}\\
G'_{\rho,\sigma}&\simeq\mathrm{Gl}'(2,\widehat{\mathbb Z})/\pm1,\quad\text{où}\quad
\mathrm{Gl}'(2,\widehat{\mathbb Z})\overset{\mathrm{dfn}}=
\{u\in\mathrm{Gl}(2,\widehat{\mathbb Z})\mid\det u\equiv1\ (3)\}\\
&\deleted{\unsure{G\,\mathbb Z\,B/\mu_2}}\\
Z'_{\rho}&=\{c\rho+d\mid c^2+cd+d^2\in\widehat{\mathbb Z}^{*}\}/\pm1\\
&\inserted{\simeq\unsure{T_\rho|\mathcal T_\rho}(\widehat{\mathbb Z})/\mu_2(\mathbb Z)}\\
&\deleted{\text{c'est le}\ \unsure{T_\rho}(\widehat{\mathbb Z})/\mu_2(\mathbb Z)\ \text{\unsure{avec notations du} \S48}}\\
Z'_{\sigma}&=\deleted{\unsure{\{t\sigma+u\mid t^2+u^2\in\widehat{\mathbb Z}^{*}\}/\pm1}\quad\text{\unsure{sous-groupe} du \S48}}\\
&\inserted{N'_{\sigma}(\widehat{\mathbb Z})/\mu_2(\mathbb Z)}\quad\text{avec notations du \S48}\\
S_0M'_{\rho,\sigma}&\simeq\deleted{\text{\unsure{est le}}}\
\unsure{SM''_{\rho,\sigma}|SM'_{\rho,\sigma}}(\widehat{\mathbb Z})/\Sigma\quad\text{avec notations du \S48}\\
M'_{\rho,\sigma}&\simeq\deleted{\unsure{SM'_{\rho,\sigma}(\widehat{\mathbb Z})}}\
M'_{\rho,\sigma}(\widehat{\mathbb Z})/\Sigma\quad\text{id}\\
\Gamma'_{\rho,\sigma}&\simeq\mathcal T'_{\rho,\sigma}(\widehat{\mathbb Z})/\Sigma_0\quad\text{avec notations du \S48}
\end{aligned}\right.
\end{aligned}
\]
\margin{c'est $G(\widehat{\mathbb Z})/\mu_2(\mathbb Z)$ avec notations du \S48}
\mspage{521}


\scanpage{523}

$3^\circ)$ Prenons $\mathcal G=\mathrm{Sl}(2,\widehat{\mathbb Z})$, avec couple modulaire.
\[
(95)\quad\left\{\begin{aligned}
B'_{\rho,\sigma}&=
\left\{\begin{pmatrix}\mu&\lambda\\0&1\end{pmatrix}\mid
\mu\in\widehat{\mathbb Z}^{*},\ \mu\equiv1\ (3),\ \lambda\in\widehat{\mathbb Z}\right\}\\
G'_{\rho,\sigma}&=\mathrm{Gl}'(2,\widehat{\mathbb Z})
=\{u\in\mathrm{Gl}(2,\widehat{\mathbb Z})\mid\det u\equiv1\ (3))\\
Z'_\rho&=\{c\rho+d\mid
\overset{\inserted{c,d\in\widehat{\mathbb Z},\ \text{t.q.}}}{c^2+cd+d^2\in\widehat{\mathbb Z}^{*}}\}\\
Z'_\sigma&=T'_\sigma\cdot\{1,\tau_\infty\},\quad\text{où}\quad
\left\{\begin{aligned}
T'_\sigma&\overset{\mathrm{dfn}}=
\{t\sigma+u\mid t,u\in\widehat{\mathbb Z},\ t^2+u^2\in\widehat{\mathbb Z}^{*}\}\\
\tau_\infty&=\begin{pmatrix}-1&0\\0&1\end{pmatrix}
\end{aligned}\right.\\
S_0M'_{\rho,\sigma}&=S_0M'_{\rho,\sigma}(\widehat{\mathbb Z})/\Sigma_0\quad\text{avec notations \S48}\\
M'_{\rho,\sigma}&=M'_{\rho,\sigma}(\widehat{\mathbb Z})/\Sigma_0\quad\rule{3cm}{0.3pt}\\
S_0\Gamma'_{\rho,\sigma}&=\unsure{S_0\mathcal T'_{\rho,\sigma}|S\mathcal T'_{\rho,\sigma}}(\widehat{\mathbb Z})/\Sigma_0\quad\rule{3cm}{0.3pt}\\
\Gamma_{\rho,\sigma}&=\mathcal T'_{\rho,\sigma}(\widehat{\mathbb Z})/\Sigma_0\quad\rule{3cm}{0.3pt}
\end{aligned}\right.
\]

\deleted{
$\overline{\text{\uline{III}}}$ \uline{Le groupe} $\unsure{M_{\rho,\sigma}|M_\rho}$

Considérons les homomorphismes

$(1)$
\diagram{\begin{tikzcd}[ampersand replacement=\&,column sep=large]
\deleted{\widehat{\mathbb Z}}
 \arrow[r,shift left=1.7em,"{\deleted{n\mapsto\rho^n}}"]
 \arrow[r,"{\deleted{n\mapsto\sigma^n}}"]
 \arrow[r,shift right=1.7em,"{\deleted{n\mapsto\epsilon_0^n}}"']
\&\deleted{\mathcal G}
\end{tikzcd}}
\deleted{\unsure{d'image} $\unsure{L_0^{\mathcal G}}$}
\diagram{\begin{tikzcd}[ampersand replacement=\&,column sep=large]
\widehat{\mathbb Z}\arrow[r]\&\mathcal G\\
n\arrow[r,mapsto]\&\epsilon_0^n
\end{tikzcd}}
où $\epsilon_0=\sigma^{-1}\rho$
(d'image $L_0^{\mathcal G}$)

et soit $J$ \deleted{l'intersection des} \inserted{\unsure{le}} noyaux, de telle façon que \deleted{\unsure{l'image} (1), \unsure{se factorise}} \deleted{\unsure{sur un}}
\[
(2)\qquad
\bigl(\widehat{\mathbb Z}^{\deleted{\unsure{*}}}/J\bigr)^*
\simeq\operatorname{Aut}(L_0^{\mathcal G}).
\]
On a \deleted{donc} un hom. évident
\[
(3)\qquad
B'_{\rho,\sigma}\longrightarrow\operatorname{Aut}L_0^{\mathcal G}
\simeq(\widehat{\mathbb Z}/J)^*
\]
nul sur $L_0$, \deleted{\unsure{i.e. se factorise}} qui donne un \unsure{hom.} correspondant
\[
(4)\qquad
\gamma'_{\rho,\sigma}\xrightarrow{\chi_\gamma}(\widehat{\mathbb Z}/J)^*,
\]
d'où en composant avec $G'_{\rho,\sigma}\longrightarrow\gamma'_{\rho,\sigma}$
\[
(4\text{bis})\qquad G'_\rho\xrightarrow{\chi_G}(\widehat{\mathbb Z}/J)^*.
\]
}
\mspage{522}


\scanpage{524}
\mspage{565}

$\overline{\text{\uline{III}}}$ \deleted{Considérons le}

Considérons l'extension $\mathrm{Gl}(2,\mathbb Z)^\wedge$ de $\mathcal G_{0,3}^{\deleted{\unsure{+}}\wedge}$ par $\{\pm1\}$:

$(1)$
\diagram{\begin{tikzcd}[ampersand replacement=\&,column sep=1.4em,row sep=1.6em]
\&\&1\arrow[d]\&1\arrow[d]\&\\
1\arrow[r]\&\{\pm1\}\arrow[r,no head]\arrow[d,equal]\&
\mathrm{Sl}(2,\widehat{\mathbb Z})\arrow[r]\arrow[d]\&
\mathcal G_{0,3}^{+\wedge}\arrow[r]\arrow[d]\&1\\
1\arrow[r,no head]\&\{\pm1\}\arrow[r,no head]\&
\mathrm{Gl}(2,\widehat{\mathbb Z})\arrow[r,no head]\arrow[d]\&
\mathcal G_{0,3}^{\wedge}\arrow[r,no head]\arrow[d]\&1\\
\&\&\{\pm1\}\arrow[r]\arrow[d]\&\{\pm1\}\arrow[d]\&\\
\&\&1\&1\&
\end{tikzcd}}

On s'intéresse au groupe $\deleted{\unsure{M}}N_{1,1}^{\sim}$, défini \inserted{extension de $\deleted{\mathcal G_{0,3}^{\wedge}}M_{0,3}^{\sim}$ par $\{\pm1\}$,} comme image inverse de l'extension $\mathrm{Gl}(2,\widehat{\mathbb Z})$ de $\mathrm{Gl}(2,\widehat{\mathbb Z})/\pm1$ par $\{\pm1\}$, via l'hom. canonique
\[
(2)\qquad
M_{0,3}^{\sim}\deleted{\mathcal G_{0,3}^{\wedge}}
\longrightarrow\mathrm{Gl}(2,\widehat{\mathbb Z})'
=\mathrm{Gl}(2,\widehat{\mathbb Z})/\pm1
\]
(cf \S45 \unsure{p.}\ldots), de sorte qu'on trouve un hom. de suites exactes,

\deleted{$\unsure{\tau_\infty(\epsilon_0)=\epsilon_0^{-1}}$}

$(3)$
\diagram{\begin{tikzcd}[ampersand replacement=\&,column sep=1.5em,row sep=2em]
1\arrow[r]\&\deleted{\unsure{M_{0,3}}}\{\pm1\}\arrow[r,no head]\arrow[d,equal]\&
\unsure{N_{1,1}^{\sim}|N_{0,3}^{\sim}}\arrow[r,no head]\arrow[d]\&M_{0,3}^{\sim}\arrow[r]\arrow[d]\&1\\
1\arrow[r]\&\{\pm1\}\arrow[r]\&\mathrm{Gl}(2,\widehat{\mathbb Z})\arrow[r]
\arrow[d,draw=none,"{\deleted{\unsure{\downarrow}}}"]\&
\mathrm{Gl}(2,\widehat{\mathbb Z})'\arrow[r,no head]
\arrow[d,draw=none,"{\deleted{\unsure{\downarrow}}}"]\&1\\
\&\&\deleted{\unsure{1}}\&\deleted{\unsure{1}}\&
\end{tikzcd}}

On trouve, \deleted{pour $\unsure{\mathcal T_{1,1}^{\wedge}}$} \inserted{\unsure{un}} comme l'image inverse de $\deleted{\mathrm{Sl}(2,\mathbb Z)}\mathcal G_{0,3}^{+\wedge}\subset$ \deleted{dans} $M_{0,3}^{\sim}$ dans $N_{1,1}$, un $\mathrm{Gl}(2,\mathbb Z)^\wedge$, dédui-
\margin{\uline{NB} $\mathcal T_{1,1}\simeq\unsure{\mathrm{Gl}(2,\mathbb Z)_{\wedge}|\uline{\mathrm{Gl}(2,\mathbb Z)}}$.}
\mspage{523}


\scanpage{525}

sant un hom
\[
(4)\qquad
\unsure{\mathrm{Sl}(2,\mathbb Z)^\wedge|\mathrm{Gl}(2,\mathbb Z)^\wedge}
\longrightarrow\deleted{\unsure{M}}N_{1,1}^{\sim}
\]
dont la restriction à $\mathrm{Sl}(2,\mathbb Z)^\wedge$ s'insère \deleted{s'insère} dans une suite exacte canonique

$(5)$
\diagram{\begin{tikzcd}[ampersand replacement=\&,column sep=1.8em]
1\arrow[r,no head]\&
\underbrace{\mathrm{Sl}(2,\mathbb Z)^{\unsure{\wedge}}}_{\mathcal T_{1,1}^{+\wedge}}\arrow[r]\&
N_{1,1}^{\sim}\arrow[r]\&
\underset{\substack{\Vert\\\Gamma_{0,3}^{\sim}}}{\Gamma_{\mathbb Q}^{\sim}}\arrow[r,no head]\&1
\end{tikzcd}}

qui à son tour s'insère dans le diagramme

$(6)$
\diagram{\begin{tikzcd}[ampersand replacement=\&,column sep=1.7em,row sep=2em]
1\arrow[r]\&\mathrm{Sl}(2,\mathbb Z)^\wedge\arrow[r]\arrow[d,"{\text{épi}}"]\&
N_{1,1}^{\sim}\arrow[r,no head]\arrow[d,"{\text{épi}}"]\&
\Gamma_{\mathbb Q}^{\sim}\arrow[r,no head]\arrow[d,"{\text{épi}}","\chi"']\&1\\
1\arrow[r]\&\mathrm{Sl}(2,\widehat{\mathbb Z})\arrow[r,no head]\&
\mathrm{Gl}(2,\widehat{\mathbb Z})\arrow[r]\&\widehat{\mathbb Z}^{*}\arrow[r,no head]\&1
\end{tikzcd}}

On peut écrire
\[
\begin{gathered}
(7)\qquad M_{0,3}^{\sim}\simeq
\overbrace{M_{0,3}^{\sim}(0)}^{\simeq\Gamma_{\mathbb Q}^{\sim}}
\cdot\overbrace{\mathcal G_{0,3}^{+\wedge}}^{\simeq\mathrm{Sl}(2,\mathbb Z)^{\wedge\prime}\simeq\mathrm{Sl}(2,\mathbb Z)^\wedge/\pm1}\\
\text{(semi-direct)}
\end{gathered}
\]
en termes de \unsure{cette} décomposition, et compte tenu qu'on a un hom canonique
\[
(8)\qquad M_{0,3}^{\sim}(0)\longrightarrow\mathrm{Gl}(2,\widehat{\mathbb Z})^{\unsure{\sharp|\prime}}
\]
\unsure{que|qui} prenant \unsure{la|là} \deleted{\unsure{hom. composé de}} \inserted{restriction de} (2) à $M_{0,3}^{\sim}(0)$, on trouve un iso. can.
\mspage{524}

\scanpage{526}
\mspage{566}

\[
 (9)\qquad
 N_{1,1}^{\sim}\simeq
 \underbrace{N_{1,1}^{\sim}(0)}_{M_{0,3}^{\sim}(0)\simeq\Gamma_{\mathbb Q}^{\sim}}
 .\,\mathrm{Sl}(2,\mathbb Z)^\wedge .
\]

Ainsi on a une opération de
\inserted{$N_{1,1}^{\sim}(0)\simeq$}$\Gamma_{\mathbb Q}^{\sim}$
sur $\mathrm{Sl}(2,\mathbb Z)^\wedge$
que je voudrais expliciter, \deleted{On va}
\unsure{utilisons|utilisant} l'\deleted{\unsure{isom.}} diagramme cartésien

\[
(10)\qquad
\diagram{\begin{tikzcd}[ampersand replacement=\&,column sep=large,row sep=large]
\mathrm{Sl}(2,\unsure{\mathbb Z|\widehat{\mathbb Z}})^\wedge
 \arrow[r] \arrow[d]
 \arrow[dr,phantom,"\text{cart}"]
 \& \mathcal{G}_{0,3}^{+\wedge}
 \arrow[r,phantom,"\simeq"]
 \arrow[dr,"{\text{induit par (2)}}"']
 \& \mathrm{Sl}(2,\mathbb Z)^\wedge/\pm1
 \arrow[d,"{\unsure{q|q'}\ \text{can}}"]
 \\
\mathrm{Sl}(2,\widehat{\mathbb Z}) \arrow[rr]
 \& {} \& \mathrm{Sl}(2,\widehat{\mathbb Z})/\pm1
\end{tikzcd}}
\]

Un élément $\unsure{\widetilde{\nu}|\widetilde{\gamma}}$ de
$\Gamma_{\mathbb Q}^{\sim}$ correspond à un couple
\deleted{$(\alpha,\beta)$ \unsure{d'élts}}

\[
(11)\qquad
(\alpha,\beta)\in
\deleted{\unsure{\mathcal{G}_{0,3}}}
\underbrace{\widehat{\Pi}_{0,3}^{\deleted{\unsure{\wedge|\prime}}\prime}}_{
 \unsure{\{1,\tau_\infty\}}.\widehat{\Pi}_{0,3}}
\times\widehat{\Pi}_{0,3}
\]

satisfaisant \deleted{\unsure{de|à}} l'équation des lacets

\[
(12)\qquad
\alpha\rho(\alpha)^{-1}=\beta\sigma(\beta)^{-1}\ell_1^\nu
\qquad (\nu\in\widehat{\mathbb Z})
\]

qui opère sur $\mathcal{G}_{0,3}^{+\wedge}\simeq
\mathrm{Sl}(2,\mathbb Z)^{\wedge\unsure{\prime}}$ par

\[
(13)\qquad
\left\{\begin{aligned}
u(\rho)&=\operatorname{int}(\alpha)\rho=\zeta\rho,
&\zeta&=\alpha\rho(\alpha)^{-1},\\
u(\sigma)&=\operatorname{int}(\beta)\sigma=\eta\sigma,
&\eta&=\beta\sigma(\beta)^{-1}
\end{aligned}\right.
\]

ce qui est compatible avec \deleted{une}\inserted{l'} opération
\deleted{(\unsure{côté} \illegible{2})}
\inserted{notée $\unsure{\underline{u}|\underline{\alpha}}$}
sur $\mathrm{Sl}(2,\widehat{\mathbb Z})$ induit
par $\operatorname{int}(\underline{u})$, où
$\underline{u}\in\mathrm{Gl}(2,\widehat{\mathbb Z})$ est l'élément

\mspage{525}

\scanpage{527}

défini \deleted{par} comme l'image de $u$ par (8), de
façon précise

\[
(14)\qquad
\unsure{\widetilde{u}|\underline{u}}=
\unsure{i_{T_0}|i_{\Gamma_0}}(\mu),\qquad
\text{où }\mu=\chi(u)\text{ est le multiplicateur, et}
\]

\[
(15)\qquad
\left\{\begin{aligned}
\unsure{i_{T_0}|i_{\Gamma_0}} &:G_m\longrightarrow
\mathrm{Gl}(2,\widehat{\mathbb Z})\\
\mu&\longmapsto\begin{pmatrix}\mu&0\\0&1\end{pmatrix}.
\end{aligned}\right.
\]

On aura, de façon précise

\[
(16)\qquad
\left\{\begin{aligned}
\underline{u}(\underline{\rho})&=\underline{\zeta}(\rho)
&\text{où }\underline{\zeta}&=\underline{\alpha}\rho(\underline{\alpha})^{-1}\\
\underline{u}(\underline{\sigma})&=\unsure{\underline{\eta}}(\sigma)
&\text{où }\underline{\eta}&=\varepsilon_{\unsure{\varphi|\rho}}
\underline{\beta}\sigma(\underline{\beta})^{-1}
\end{aligned}\right.
\]

\margin{les soulignés désignent les éléments de
$\mathrm{Gl}(2,\widehat{\mathbb Z})$ ---
$\underline{\alpha}$ et $\underline{\beta}$
\unsure{ne sont déterminés} \unsure{qu'au|qu'à un} signe près,
mais $\underline{\zeta},\underline{\eta}$ sont
\unsure{sans ambiguïté|univoques}.}

\deleted{$\varepsilon=\varepsilon_2(\mu)$}

avec, pour mémoire

\[
(17)\qquad
\varepsilon=\varepsilon_2(\mu)\in\pm1
\qquad
\left\{\begin{aligned}
\varepsilon&=+1&&\text{ssi }\mu\equiv1\ (4)\\
\varepsilon&=-1&&\text{ssi }\mu\equiv-1\ (4)
\end{aligned}\right.
\]

Bien entendu, $\mathrm{Sl}(2,\mathbb Z)^\wedge$ est
\uline{engendré} \inserted{\unsure{top.}} par $\rho$ et $\sigma$,
et on identifie $\pm1$ au sous-groupe \unsure{central|centre}
de $\mathrm{Sl}(2,\mathbb Z)$ donc de
$\mathrm{Sl}(2,\mathbb Z)^\wedge$, on écrira aussi
$-g$ pour le produit de $g\in\mathrm{Sl}(2,\mathbb Z)^\wedge$
par $-1$. On aura donc dans $\mathrm{Sl}(2,\mathbb Z)^\wedge$

\[
(18)\qquad
\left\{\begin{aligned}
u(\rho)&=\deleted{\operatorname{int}\unsure{\zeta}(\rho)=}
\operatorname{int}(\widetilde{\alpha})(\rho)=\zeta\rho\\
u(\unsure{\sigma|\rho})&=
\varepsilon\operatorname{int}(\widetilde{\beta})(\sigma)=\eta\sigma
\end{aligned}\right.
\]
\unsure{et}

avec
\[
(19)\qquad
\left\{\begin{aligned}
\zeta&=\widetilde{\alpha}\rho(\widetilde{\alpha})^{-1},\qquad
\eta=\varepsilon\widetilde{\beta}\sigma(\widetilde{\beta})^{-1}\\
&\varepsilon=\varepsilon_2(\mu)\in\{\pm1\}\subset
\mathrm{Sl}(2,\mathbb Z)^\wedge,
\end{aligned}\right.
\]

\mspage{526}

\scanpage{528}
\mspage{567}

où maintenant $\widetilde{\alpha}$ et $\widetilde{\beta}$
désignent des relèvements arbitraires de $\alpha$ et $\beta$ à
$\mathrm{Sl}(2,\mathbb Z)^\wedge$ --- ils ne sont déterminés
qu'au signe près, mais
\inserted{$\operatorname{int}(\widetilde{\alpha}),
\operatorname{int}(\widetilde{\beta}),$}
$\zeta$ et $\eta$ sont déterminés sans ambiguïté par $u$.
\deleted{Si \unsure{maintenant}}
\inserted{Notons aussi:}
\[
(20)\qquad u(\varepsilon_0)=\varepsilon_0^\mu
\quad\deleted{\text{\unsure{car}}}
\]

Si maintenant on a un hom \inserted{\unsure{surj}}
\[
(20)\qquad N_{1,1}^{\sim}\xrightarrow{\Psi}G
\]
induisant un hom
\[
(21)\qquad
\mathrm{Sl}(2,\mathbb Z)^\wedge
\xrightarrow{\text{épi}}
\underset{\substack{\Vert\\\operatorname{Im}\Psi}}{\mathcal{G}}
\subset G,
\]

les images de $\rho,\sigma$ dans $\mathcal{G}$
(notées par les $\overline{\mathrm m}$ lettres)
\unsure{satisfaisant|satisfont}
\[
(22)\qquad
\rho^6=\sigma^4=1,\qquad
\rho^3=\sigma^2\quad\text{et par suite}
\]
\deleted{$[\rho^3,\sigma^2]=1$ \unsure{ce qui suffit}}
\[
[\rho^3,\sigma]=[\sigma^2,\rho]=1
\quad\text{i.e. }\rho^3,\sigma^2
\]
\deleted{à \unsure{impliq.}}
sont \uline{centraux}

\margin{donc dans $\mathcal{G}$ il y a un élément central
canonique involutif, noté $-1$ \unsure{. .}}

et on peut, pour $u\in N_{1,1}^{\sim}(0)$, correspondant à des
\unsure{élmts} $\widetilde{\alpha},\widetilde{\beta}$ de
$\mathrm{Sl}(2,\mathbb Z)^\wedge$
(déterminés mod \uline{signe}), d'où $\zeta,\eta$,
essayer de trouver un \inserted{\unsure{endom.|autom.}}
$u_{\mathcal{G}}$ de $\mathcal{G}$,
satisfaisant les relations idoines
\[
(23)\qquad
u_{\mathcal{G}}(\rho_{\mathcal{G}})
=\zeta_{\mathcal{G}}\rho_{\mathcal{G}},\qquad
u_{\mathcal{G}}(\sigma_{\mathcal{G}})
=\eta_{\mathcal{G}}\sigma_{\mathcal{G}},
\]
et faire jouer \unsure{nos|aux} $\widetilde{\alpha}$ et
$\widetilde{\beta}$. Mais on est

\mspage{527}

\scanpage{529}

obligé, comme \deleted{\unsure{en général}} on n'a pas
\inserted{\unsure{néc\textsuperscript{t}.}}
$\widetilde{\alpha}\in\mathrm{Sl}(2,\mathbb Z)^\wedge$,
seulement $\widetilde{\alpha}\in\mathrm{Gl}(2,\mathbb Z)^\wedge$,
de faire intervenir l'élément $\tau_\infty$ de
$\mathrm{Gl}(2,\mathbb Z)^\wedge$,
et \inserted{l'action de} son image $\tau$ dans $G$ sur
$\mathcal{G}$. On aura, en prenant maintenant

\[
(24)\qquad
\begin{aligned}
\overset{\deleted{\unsure{\mathcal{G}|\wedge}}}{\varepsilon_{0\mathcal{G}}}
&=\overset{\deleted{\unsure{\mathcal{G}|\wedge}}}{\sigma_{\mathcal{G}}}
\overset{\deleted{\unsure{\mathcal{G}|\wedge}}}{\rho_{\mathcal{G}}},
&\varepsilon_{1\mathcal{G}}&=\rho_{\mathcal{G}}\sigma_{\mathcal{G}}
\deleted{\unsure{=}}\,=\sigma_{\mathcal{G}}(\varepsilon_{0\mathcal{G}})
=\rho_{\mathcal{G}}(\varepsilon_{0\mathcal{G}})\\
&=-\sigma_{\mathcal{G}}^{-1}\rho_{\mathcal{G}}
&&=-\rho_{\mathcal{G}}\sigma_{\mathcal{G}}^{-1}
\end{aligned}
\]

\margin{ce qui n'est pas la convention habituelle
\unsure{jusqu'à} présent\ldots}

les relations
\[
(25)\qquad
\tau(\sigma)=\sigma^{-1},\qquad
\tau(\rho)=\sigma(\rho^{-1})
\]

\margin{je laisse tomber les indices $\mathcal{G}$}

qui impliquent, via (24) et (22)
\[
(26)\qquad
\tau(\varepsilon_0)=\varepsilon_0^{-1},\qquad
\tau(\varepsilon_1)=\varepsilon_1^{-1},
\]
et on introduit le groupe
\[
(27)\qquad
\deleted{\unsure{\{1,\tau\}}}
\mathcal{G}^{\unsure{\bullet|\prime}}
\overset{\mathrm{dfn}}{=}
\underset{\frac12\,\text{direct}}{\{1,\tau\}.\mathcal{G}},
\]
où $\{1,\tau\}$ opère sur $\mathcal{G}$ par les formules (25).

Ceci dit, on est en mesure de paraphraser des constructions
déjà faites par ailleurs.

\mspage{528}
